% Generated draft: requires mathematical and source review.
\documentclass[11pt]{article}
\usepackage[margin=25mm]{geometry}
\usepackage{fontspec}
\setmainfont{Latin Modern Roman}
\usepackage{amsmath,amssymb,mathtools,unicode-math}
\setmathfont{Latin Modern Math}
\usepackage{enumitem,needspace,xurl,hyperref}
\setlength{\parindent}{0pt}
\setlength{\parskip}{5pt}
\setlength{\emergencystretch}{3em}
\newcommand{\N}{\mathbb{N}}
\newcommand{\Z}{\mathbb{Z}}
\newcommand{\Q}{\mathbb{Q}}
\newcommand{\R}{\mathbb{R}}
\newcommand{\C}{\mathbb{C}}
\newcommand{\F}{\mathbb{F}}
\newcommand{\A}{\mathbb{A}}
\newcommand{\PP}{\mathbb{P}}
\newcommand{\E}{\mathbb{E}}
\newcommand{\eps}{\varepsilon}
\newcommand{\dd}{\,\mathrm{d}}
\begin{document}
\section*{UnsolvedMath ID: 30005339}
\textbf{Source problem code:} catalogue entry
\subsection{Definitions and mathematical statement}
The supplied record concerns the number of real zeros of random functions, with random polynomials mentioned in its background. It does not specify a family of functions, their probability laws, or their domains. Consequently, it supplies a research programme rather than a fully determined mathematical problem.

To express the target conditionally, let \(f_n\) be a specified family of random functions on specified sets \(D_n\subseteq\R\). If distinct zeros are intended, define
\[
N_n=\#\{x\in D_n:f_n(x)=0\}.
\]
This is a possible counting convention, not a convention established by the supplied record. Counting zeros with multiplicity would require a different definition and sufficient regularity to define multiplicity. Measurability and almost-sure finiteness of the count must also be established for the chosen model.

\textbf{Conditional mathematical target.} For each intended model and counting convention, determine the expectation \(\E[N_n]\), the variance
\[
\operatorname{Var}(N_n)=\E\bigl[(N_n-\E[N_n])^2\bigr],
\]
when these quantities are defined, and their asymptotic behaviour in a specified limiting regime. Determine suitable deterministic centering constants \(a_n\in\R\), scales \(b_n>0\), and a probability law \(L\), if such choices exist, for which
\[
\frac{N_n-a_n}{b_n}\xrightarrow{\mathrm d}L
\qquad(n\to\infty).
\]
The sequential indexing, limiting regime, and affine normalization above are an explanatory framework; the record does not specify them or assert that such a limit exists. In particular, no Gaussian limit is stipulated.

\textbf{Unresolved source qualifications.} The original excerpt refers to \(N_{n,\xi}\) without defining either the random function or \(\xi\). It contains an OCR-corrupted logarithmic bound and an expectation formula attributed to Kac, involving integrals and \(p(t,0,y)\). The bound cannot be recovered reliably from the supplied text, and the meaning of \(p\) and the hypotheses of the formula are absent. Neither is asserted here as a mathematical result. The intended classes of functions, coefficient assumptions, zero-counting convention, observation domain, and asymptotic parameter remain unspecified.
\subsection{Short English statement}
For specified random-function models, determine the mean, variance, and suitably normalized limiting distribution of their real-zero counts. The supplied record leaves the models, counting convention, and limiting regime unspecified, so this is a conditional formulation requiring source clarification and human mathematical review.
\subsection{Sources}
\begin{itemize}
\item Local problem registry: \texttt{data/unsolvedmath/problems.json}, with \texttt{catalog/problems.json} as the additional registry specified by the project; supplied complete record for permanent UnsolvedMath ID \(30005339\), title ``Real Zeros of General Random Functions,'' exact source problem code \texttt{catalogue entry}. No individual proposer is supplied; the recorded proposed year is 2023.
\item Supplied report citation: \emph{Combinatorics} (2023), Oberwolfach Reports, DOI 10.4171/owr/2023/1, pp.\ 5--89. \url{https://doi.org/10.4171/owr/2023/1}. The record describes the report as documenting workshop discussions and open problems; no problem-specific page or speaker attribution is provided.
\item Supplied literature source: Do--Nguyen, \emph{Real Roots of Random Polynomials}. \url{https://arxiv.org/abs/2303.05478}.
\item Supplied literature source: \emph{Random Polynomial Central Limit Theorems}. \url{https://arxiv.org/abs/1904.04347}. No author attribution is supplied for this reference.
\item Literature assessment dated exactly \texttt{2026-08-22}: the supplied curation labels the problem partially solved and states that classical and several structured coefficient models have asymptotic means, variances, and central limit theorems, while the general random-polynomial or random-function question has no complete solution. This is supplied dated curation, not an independently verified literature assessment. The background explicitly requests independent verification of time-sensitive status and cited claims.
\item The original excerpt attributes early work on this question to Kac and Littlewood--Offord in the 1940s, associates Littlewood--Offord's anti-concentration inequality with a proof discussed there, and attributes a general expectation formula to Kac. These historical claims are supplied attribution; no specific original-paper citations or complete theorem hypotheses are included.
\item The registry marks the statement \texttt{corrected\_verified} and reports that it was checked against the official report and preserves its hypotheses, notation, and target. This is a supplied verification claim, not verification performed for this draft; the supplied excerpt does not contain enough information to recover those hypotheses.
\end{itemize}
\end{document}
